
\section{Introduction}

The software engineering landscape is currently undergoing a significant change due to the rapid evolution of Artificial Intelligence. Specifically, the rise of Generative AI and Large Language Models has extended beyond simple classification tasks to enable the generation of coherent, context-aware outputs, reshaping traditional problem-solving and creative workflows across the technology sector.

Within this context, these technologies are revolutionizing the Software Development Life Cycle (SDLC) by introducing "Vibe Coding" and intent-driven programming, where developers collaborate with models using natural language. This paradigm shift automates complex tasks from code generation to documentation, but also necessitates a transition from manual code construction to a role focused on oversight, validation, and architectural decision-making to maintain high standards of reliability, security and quality.

To investigate the Software Development environment landscape rigorously, this report employs a dual methodological approach comprising a Systematic Literature Review (SLR) and the Design Science Research Methodology (DSRM). The SLR, following Kitchenham’s guidelines, was conducted to map the current state of AI adoption, benefits, and consequences, while the DSRM provides the structured process for designing, developing, and validating the proposed solution artifact.

Despite the rapid adoption of these tools, the industry faces a critical problem, and therefore it is defined a research problem: the absence of established, solid, and clear methodologies to guide the responsible integration of AI and LLMs into professional workflows. Existing research remains fragmented across isolated tools, leaving organizations without a structured approach to manage the risks regarding trust, explainability, and long-term maintainability that accompany AI-assisted development.

Addressing this gap, this report proposes a new framework inspired by Agile and Scrum principles designed to support the adoption, collaboration, and governance of AI-assisted development practices. This artifact defines specific phases, practices, and oversight mechanisms to integrate AI throughout the SDLC, ensuring that the technology enhances productivity without compromising value or accountability.

The applicability of this framework is demonstrated through a realistic project scenario within a next-generation consultancy. This demonstration illustrates how the proposed artifact functions in a real-world setting, guiding a team through the challenges of Vibe Coding and AI-augmented prototyping while maintaining alignment with business goals.

To validate the proposal, an ex post evaluation strategy is outlined, involving interviews with practitioners and Agile experts alongside an analytical assessment against defined success criteria. This evaluation aims to confirm the framework's utility, feasibility, and alignment with industry needs, ensuring it effectively supports Human-AI collaboration.

The remainder of this report is structured to guide the reader through the research process: Section 2 details the methodologies used; Section 3 provides the theoretical background; Section 4 presents the findings of the Systematic Literature Review; Section 5 defines the research problem; Section 6 outlines the objectives; Section 7 introduces the proposed framework; and Sections 8 through 11 cover the demonstration, evaluation plan, communication strategy, and final conclusions.



